\chapter{Áp Lực Đồng Trang Lứa}
\markboth{Áp Lực Đồng Trang Lứa}{Áp Lực Đồng Trang Lứa}

\section{Bạn Em Giỏi Quá Nên Em Nản}
\begin{truyen}{Bạn Em Giỏi Quá Nên Em Nản}{My Friends Are Too Good, So I Gave Up}
\chuhoa{H}{ọc sinh:} ``Nhìn mấy bạn giỏi quá em nản, thầy ạ. Làm gì cũng thấy mình thua từ đầu.''
\textit{(Student: ``Seeing my classmates so good makes me discouraged. Whatever I do, I already feel behind.'')}

\teacherpair{Giỏi của người khác có thể làm em áp lực. Nhưng nếu em lấy nó làm lý do để buông luôn phần của mình, thì em đang biến sự xuất sắc của người khác thành công cụ phá chính mình. Không ai bắt em phải nở cùng một tốc độ. Nhưng em cũng không được quyền nằm luôn rồi gọi đó là bất công.}{``Other people's excellence can pressure you. But if you turn that into a reason to abandon your own effort, then you are using their success as a tool to damage yourself. No one requires you to bloom at the same speed. But you also do not have the right to lie down and call that injustice.''}
\end{truyen}

\section{Đồng Trang Lứa Hay Đồng Sân Khấu}
\begin{truyen}{Đồng Trang Lứa Hay Đồng Sân Khấu}{Peers or an Audience?}
\chuhoa{H}{ọc sinh:} ``Đi học giờ giống lên sân khấu hơn là đi học. Ai cũng nhìn nhau, so nhau, khoe nhau.''
\textit{(Student: ``School now feels more like a stage than a classroom. Everyone watches, compares, and shows off.'')}

\teacherpair{Em nhìn ra như vậy là đúng. Nhưng càng ở chỗ ai cũng diễn, em càng phải biết mình học vì cái gì. Nếu em để toàn bộ giá trị bản thân bị buộc vào ánh nhìn của bạn bè, em sẽ sống rất mệt và rất nông. Vì đám đông tuổi học trò thay gu nhanh hơn cả thời tiết.}{``You are right to notice that. But the more you are surrounded by performance, the more you must know what you are learning for. If your entire self-worth depends on your peers' gaze, you will live a shallow and exhausting life. Student crowds change taste faster than the weather.''}
\end{truyen}

\section{Áp Lực Có Thể Đẩy, Không Nhất Thiết Phải Đè}
\begin{truyen}{Áp Lực Có Thể Đẩy, Không Nhất Thiết Phải Đè}{Pressure Can Push, Not Only Crush}
\chuhoa{H}{ọc sinh:} ``Vậy em phải làm sao với áp lực đồng trang lứa?''
\textit{(Student: ``Then what should I do with peer pressure?'')}

\teacherpair{Chọn vài người đáng học, không phải vài người đáng ghen. So với chính mình của tháng trước nhiều hơn so với bài đăng của người khác hôm nay. Và nhớ rằng: áp lực sẽ đè em chết nếu em chỉ đứng dưới nó. Nhưng nó có thể đẩy em đi nếu em biết đứng nghiêng đúng hướng.}{``Choose a few people worth learning from, not a few people worth envying. Compare yourself more with who you were last month than with someone else's post today. Pressure crushes you if you only stand beneath it. But it can also push you forward if you lean in the right direction.''}
\end{truyen}

\begin{baihoc}
Bạn giỏi không phải là cái cớ để em bỏ cuộc.
\textit{(Other people's talent is not an excuse for your surrender.)}

Áp lực đồng trang lứa nguy hiểm nhất khi nó cướp mất nhịp phát triển riêng của mình.
\textit{(Peer pressure becomes most dangerous when it steals your own rhythm of growth.)}
\end{baihoc}

\begin{ghinhoanh}
\textbf{Short English to remember:}

Comparison can guide.

It should not define you.
\end{ghinhoanh}

\begin{tuhocnhanh}
1. Em đang nhìn ai với tâm thế học hỏi, và đang nhìn ai với tâm thế ghen tị? \textit{(Whom do you look at to learn, and whom do you look at with envy?)}

2. Nhịp tiến bộ riêng của em đã bị đám đông làm méo đến mức nào? \textit{(How much has the crowd distorted your own pace of progress?)}
\end{tuhocnhanh}

\ngancach
